\needspace{4\baselineskip}
\section{Benchmark Details}\label{sec:appendix-benchmark}

\subsection{Code Search Checkpoint Specifications}\label{sec:appendix-code-search}

This subsection provides the complete specifications for the first three checkpoints of the \texttt{code\_search} problem referenced in \autoref{sec:running-example}. Each specification appears exactly as presented to agents during evaluation.

\subsubsection{Checkpoint 1: Python-only Exact and Regex Matching}
\label{sec:appendix-cp1}

\begin{lstlisting}[style=specblock, caption={Specification for \texttt{code\_search} checkpoint 1}, label={lst:code-search-cp1}]
Build a command-line code searcher for Python codebases.
It takes a directory of source files and a set of rules,
then prints one JSON object per match (JSON Lines) to STDOUT.

Supported language (by file extension only): Python (.py).
Supported rule types (exactly these two): exact match and
vanilla regex.

# Deliverable

Write an executable that accepts:

  <entry_command> <root_dir> --rules <rules_file> [--encoding <name>]

* <root_dir>: path to the codebase to scan.
* --rules <rules_file>: path to a JSON array of rules (see schema).
* --encoding <name>: optional; default utf-8. Files that fail to
  decode are skipped.

Output: JSON Lines to STDOUT, one object per match (schema below).
On success (even if zero matches): exit code 0.

## Inputs

### File type -> language

Only files with extension .py are scanned. All other files are
ignored.

### Rules file (JSON array)

Each rule is an object with this schema:

{
  "id": "<non-empty string>",
  "kind": "exact" | "regex",
  "pattern": "<non-empty string>",
  "languages": ["python"],                 // optional; default: ["python"]
  "regex_flags": ["i", "m", "s"]           // optional; only for kind = "regex"
}

Constraints on types/inputs

* id: unique across the rules array.
* pattern: valid UTF-8 string. For kind="regex", it must compile
  with provided regex_flags (subset of i case-insensitive, m multiline,
  s dotall). No other flags allowed.
* languages (when present) must be an array of strings and may only
  contain "python".
* Missing optional fields use their defaults.

### Matching semantics

Matches inside comments/strings are allowed.

## Output (JSON Lines to STDOUT)

For each match, print a single JSON object with exactly:

{
  "rule_id": "<string>",
  "file": "<posix path>",                 // path relative to root_dir, '/' separators
  "language": "python",
  "start": {"line": <int>, "col": <int>}, // 1-based line/column.
  "end": {"line": <int>, "col": <int>},   // position immediately AFTER the match
  "match": "<string>"
}

No extra fields. Each object on its own line with a trailing newline.
No other STDOUT output.

## Normalization

1. Ordering: Matches must appear by file (lexicographically), then
   start.line, then start.col, then rule_id.
2. Coordinates: Lines and columns are 1-based.
3. Path format: file is relative to <root_dir> with '/' separators.
4. Encoding: Read files using --encoding (default utf-8); skip files
   that fail to decode.
\end{lstlisting}

\subsubsection{Checkpoint 2: Multi-language Support with Filtering}
\label{sec:appendix-cp2}

\begin{lstlisting}[style=specblock, caption={Specification for \texttt{code\_search} checkpoint 2}, label={lst:code-search-cp2}]
Extend your code searcher to support JavaScript and C++
source files.

---

## New Requirements

### File type -> language

Scan these extensions:

| Language   | Extensions                                   |
| ---------- | -------------------------------------------- |
| Python     | .py                                          |
| JavaScript | .js, .mjs, .cjs                              |
| C++        | .cc, .cpp, .cxx, .hh, .hpp, .hxx             |

### Rule language filtering

Rules may specify "languages" from this set:
["python", "javascript", "cpp"].
If omitted, the rule applies to all three.

---

## Example

### rules.json

[
  {"id":"todo","kind":"exact","pattern":"TODO:"},
  {"id":"printf","kind":"regex","pattern":"\\bprintf\\s*\\(","languages":["cpp"]},
  {"id":"console-log","kind":"regex","pattern":"console\\.log\\s*\\(","languages":["javascript"]}
]

### Project tree

repo/
  main.py
  app.js
  src/
    engine.cpp

### Run

$ <entry_command> repo --rules rules.json

### Output

{"rule_id":"console-log","file":"app.js","language":"javascript","start":{"line":5,"col":3},"end":{"line":5,"col":15},"match":"console.log("}
{"rule_id":"printf","file":"src/engine.cpp","language":"cpp","start":{"line":42,"col":3},"end":{"line":42,"col":10},"match":"printf("}
{"rule_id":"todo","file":"main.py","language":"python","start":{"line":1,"col":1},"end":{"line":1,"col":6},"match":"TODO:"}
\end{lstlisting}

\subsubsection{Checkpoint 3: Structure-Aware Metavariable Patterns}
\label{sec:appendix-cp3}

\begin{lstlisting}[style=specblock, caption={Specification for \texttt{code\_search} checkpoint 3}, label={lst:code-search-cp3}]
Extend your code searcher to support structure-aware patterns
with metavariables.

Supported rule kinds: exact, regex, and pattern.

# Deliverable

Your existing executable is extended to understand kind: "pattern"
in the rules file:

  <entry_command> <root_dir> --rules <rules_file> [--encoding <name>]

Output: JSON Lines (one object per match).

## Pattern Rules

### Rule schema additions

{
  "id": "<non-empty string>",
  "kind": "pattern",
  "pattern": "<code-like string with metavariables>",
  "languages": ["cpp" | "javascript" | "python"],      // optional; default: all 3
}

### Metavariables

* Any token of the form $NAME (e.g., $X, $GREETING) is a metavariable.
  * Metavariables ending in ? (e.g., $X?) are optional -- the pattern
    matches even if they are not present.
* A metavariable matches a single code element appropriate for that
  position.
* If the same metavariable name appears multiple times in the pattern,
  all occurrences must match the same text.
* $$ in the pattern matches a literal $ in the source code.

### Pattern string

* The pattern must be valid code in the target language (with
  metavariables treated as valid placeholders).
* No wildcards beyond metavariables are required (no ellipsis or
  quantifiers).

### Matching semantics

Find all matches in source files where the rule's language applies.

## Output (JSON Lines to STDOUT)

For kind: "pattern", the JSON object per match is:

{
  "rule_id": "<string>",
  "file": "<posix path>",
  "language": "cpp" | "javascript" | "python",
  "start": {"line": <int>, "col": <int>},
  "end": {"line": <int>, "col": <int>},
  "match": "<string>",
  "captures": {
    "$NAME": {
      "text": "<matched source text>",
      "ranges": [
        {"start": {"line": <int>, "col": <int>}, "end": {"line": <int>, "col": <int>}}
      ]
    }
    // ... one entry per metavariable bound in this match
  }
}

* ranges lists every occurrence of that metavariable within the matched
  region (use the same 1-based, Unicode column coordinates).
* If a metavariable appears once, ranges has a single range.

## Normalization

1. Pattern determinism: When multiple matches share the same start
   position, sort by end position (earlier end first), then by rule_id.
2. Captures key order: Serialize captures with keys sorted
   lexicographically by metavariable name (e.g., $A before $X).
\end{lstlisting}
\newpage
% Generated problem design-decision overview.
% Include with \input{sections/generated/problem_design_decisions_table}

\begin{table*}[t]
\centering
\caption{SlopCodeBench problems and the design decisions stressed by iterative checkpoints, part 1 of 2. \textbf{CPs} gives the number of checkpoints.}
\label{tab:problem-design-decisions-a}
\scriptsize
\setlength{\tabcolsep}{4pt}
\renewcommand{\arraystretch}{1.08}
\begin{tabularx}{\textwidth}{@{}l c >{\raggedright\arraybackslash}X@{}}
\toprule
Problem & CPs & Design decisions stressed \\
\midrule
\texttt{cfgpipe} & 6 & Source-priority resolution, typed parsing boundaries, nested schema representation, watch and event architecture, redaction, and store-prefix composition. \\
\texttt{circuit\_eval} & 8 & Circuit IR, parser separation across formats, two-valued versus three-valued logic, optimization pass pipeline, and equivalence checking. \\
\texttt{code\_search} & 5 & Regex versus AST matching, metavariable representation, rewrite conflict handling, selector semantics, and language adapter boundaries. \\
\texttt{dag\_execution} & 3 & DAG DSL parsing, task parameter model, execution ordering, cache key design, and per-task cache overrides. \\
\texttt{database\_migration} & 5 & Migration operation model, rollback representation, dependency sorting, cycle errors, and constraint or index handling. \\
\texttt{datagate} & 7 & Dataset ingestion abstraction, query and filter pipeline, pagination and export contracts, cache invalidation, runtime config, and access control. \\
\texttt{dynamic\_buffer} & 4 & Example-to-transform inference, generated API shape, stateful transform classes, streaming buffers, and cache behavior. \\
\texttt{dynamic\_config\_service\_api} & 4 & Versioned config storage, inheritance and deep merge rules, schema validation, approval workflow state, and policy enforcement. \\
\texttt{env\_manager} & 5 & Declarative provisioning IR, OS-specific planning boundaries, deterministic output, module validation, profiles, and build-script generation. \\
\texttt{etl\_pipeline} & 5 & Pipeline IR, expression language design, branching, reusable sub-pipeline parameters, and library loading. \\
\texttt{eve\_industry} & 6 & SDE normalization, recipe graph representation, invention probability, ME and TE accounting, job scheduling, and recursive build planning. \\
\texttt{eve\_jump\_planner} & 3 & Route graph modeling, fuel and fatigue costs, spatial distance calculations, cloak mechanics, and high-sec fallback routing. \\
\texttt{eve\_market\_tools} & 4 & Market order ingestion, price-book aggregation, reprocessing yield math, hub optimization, and arbitrage or hauling profit ranking. \\
\texttt{eve\_route\_planner} & 3 & Warp physics, route cost calculation, cargo manifest planning, multi-trip scheduling, and contract selection under constraints. \\
\texttt{execution\_server} & 6 & Command execution boundaries, file tracking, output format abstraction, chain and hook semantics, cache keys, persistent environment concurrency, and job queue DAGs. \\
\texttt{file\_backup} & 4 & Schedule schema, backup strategy abstraction, glob exclusion semantics, verification, incremental hashing, and destination layout. \\
\texttt{file\_merger} & 4 & Multi-format reader abstraction, schema alignment, external sort strategy, partition and shard layout, and nested type flattening. \\
\texttt{file\_query\_tool} & 5 & SQL parsing and execution planning, file table abstraction, joins, aggregation, window state, CTEs, and subquery scope. \\
\bottomrule
\end{tabularx}
\end{table*}

\begin{table*}[t]
\centering
\caption{SlopCodeBench problems and the design decisions stressed by iterative checkpoints, part 2 of 2.}
\label{tab:problem-design-decisions-b}
\scriptsize
\setlength{\tabcolsep}{4pt}
\renewcommand{\arraystretch}{1.08}
\begin{tabularx}{\textwidth}{@{}l c >{\raggedright\arraybackslash}X@{}}
\toprule
Problem & CPs & Design decisions stressed \\
\midrule
\texttt{forge} & 8 & Resource models for blueprints, allocations, units, modules, revision gates, strict JSON contracts, and admin access boundaries. \\
\texttt{l2m} & 5 & LaTeX parsing depth, macro expansion, environment conversion rules, CLI flags, and KaTeX-compatible Markdown output. \\
\texttt{layered\_config\_synthesizer} & 4 & Layer merge semantics, conflict resolution, fragment expansion, environment interpolation, multi-run manifests, and JSON Schema validation. \\
\texttt{log\_query} & 5 & SQL-like grammar design, NDJSON streaming execution, aggregation and joins, schema label mapping, and subquery scope. \\
\texttt{meshctl} & 8 & Declarative resource specs, persisted state shape, lifecycle and topology operations, security model, migrations, network exposure, and multi-region policy. \\
\texttt{metric\_transform\_lang} & 5 & DSL parser and interpreter design, stream processing, aggregation and window state, temporal joins, and resumable execution. \\
\texttt{migrate\_configs} & 5 & Rule engine design, multi-format parsing, inheritance and cycle detection, array pattern matching, relocation rules, and validation order. \\
\texttt{mocked\_http} & 8 & Mock schema design, request matching, JSONPath evaluation, environment interpolation, validation strategy, and admin runtime behavior. \\
\texttt{mvvault} & 6 & Versioned catalog schema handling, field-level history, sync and download boundaries, viewer routing, annotations, and auto-migration triggers. \\
\texttt{pwd\_manager} & 5 & Vault storage and encryption boundary, unlock lifecycle, search and edit model, categories, clipboard integration, import or export, and locking behavior. \\
\texttt{recli} & 8 & CLI command tree architecture, config inheritance, aliases, output formatting, file cache and SQLite persistence, container orchestration, upgrades, and system checks. \\
\texttt{rejector} & 5 & Rate-limit scheduler, task and generation scheme abstractions, ICL setup representation, agentic tool-loop state, and multi-provider routing. \\
\texttt{sheeteval} & 7 & Rule schema for grading, formula and dependency evaluation, fatal and concealed controls, threshold semantics, scenario check mode, and HTML reports. \\
\texttt{sith} & 6 & Static versus interpreter-assisted analysis, symbol indexing, refactor safety, project root discovery, and environment or settings handling. \\
\texttt{test\_translator} & 8 & Structured test spec parsing, generator per target language, test discovery, equality semantics, generate-before-test enforcement, and loop-as-test translation. \\
\texttt{textdrop} & 6 & HTTP boundary validation, Markdown rendering, metadata and TOC extraction, static assets, auth and drain lifecycle, and storage backend abstraction. \\
\texttt{trajectory\_api} & 5 & Trajectory data model, token and cost tracking, mutable updates with ETags, fork lineage, grammar parsing for tool calls, and sandboxed execution. \\
\texttt{xjq} & 5 & Multi-format input detection, XPath and CSS selector dispatch, text extraction modes, file precedence, and XML or JSON output formatting. \\
\midrule
\multicolumn{1}{@{}l}{\textbf{Total}} & \textbf{196} & \textbf{36 problems} \\
\bottomrule
\end{tabularx}
\end{table*}

